\documentclass[11pt]{article}

\usepackage{amsmath,amssymb}
\usepackage[a4paper,margin=1in]{geometry}

% Remove paragraph indentation
\setlength{\parindent}{0pt}

\begin{document}

\title{A Concise Derivation of Principal Component Analysis}
\author{}
\date{}
\maketitle

\section*{Goal}

Given data points
\[
x_i \in \mathbb{R}^p,
\qquad i=1,\dots,n,
\]
PCA aims to construct a lower-dimensional representation
\[
y_i = U_k^T(x_i-m),
\qquad
y_i \in \mathbb{R}^k,
\qquad
U_k \in \mathbb{R}^{p\times k},
\qquad
k < p.
\]

Here, \(m\) is the mean of the data, and the columns of \(U_k\) define the \(k\)-dimensional subspace onto which the data is projected.

The goal is to choose \(U_k\) such that the projection preserves as much information as possible.

In PCA, information is measured by variance. Therefore, PCA chooses the projection directions so that the projected data has maximal variance.

\section*{Step 1: Center the Data}

The mean of the data is
\[
m = \frac{1}{n}\sum_{i=1}^n x_i.
\]

Define the centered data points
\[
z_i = x_i - m.
\]

Centering is necessary because PCA should describe how the data varies around its mean, rather than where the whole dataset is located in space.

Thus, the PCA projection becomes
\[
y_i = U_k^T z_i.
\]

\section*{Step 2: Start with a One-Dimensional Projection}

First consider projecting the data onto one direction
\[
v \in \mathbb{R}^p.
\]

We require
\[
\|v\|=1,
\]
or equivalently
\[
v^T v = 1.
\]

This constraint is necessary because otherwise the projected variance could be made arbitrarily large simply by scaling \(v\).

The projection of \(z_i\) onto \(v\) is
\[
v^T z_i.
\]

Therefore, the variance of the projected data is
\[
\operatorname{Var}(v)
=
\frac{1}{n-1}
\sum_{i=1}^n (v^T z_i)^2.
\]

\section*{Step 3: Write the Objective in Matrix Form}

Define the centered data matrix
\[
Z = [z_1,z_2,\dots,z_n]
\in \mathbb{R}^{p\times n}.
\]

Define the scatter matrix
\[
S = ZZ^T.
\]

Then
\[
\sum_{i=1}^n (v^Tz_i)^2
=
v^TZZ^Tv
=
v^TSv.
\]

Hence,
\[
\operatorname{Var}(v)
=
\frac{1}{n-1}v^TSv.
\]

Since \(\frac{1}{n-1}\) is constant, maximizing the variance is equivalent to solving
\[
\max_v \quad v^TSv
\]
subject to
\[
v^Tv=1.
\]

\section*{Step 4: Solve the Constrained Optimization Problem}

This is a constrained optimization problem, so we use the method of Lagrange multipliers.

Define the Lagrangian
\[
\mathcal{L}(v,\lambda)
=
v^TSv-\lambda(v^Tv-1).
\]

To find an optimum, take the derivative with respect to \(v\) and set it to zero:
\[
\nabla_v \mathcal{L}=0.
\]

Since \(S\) is symmetric,
\[
\nabla_v(v^TSv)=2Sv,
\]
and
\[
\nabla_v(v^Tv)=2v.
\]

Therefore,
\[
2Sv-2\lambda v=0.
\]

So
\[
Sv=\lambda v.
\]

This is an eigenvalue equation. Hence, the optimal projection directions must be eigenvectors of the scatter matrix \(S\).

\section*{Step 5: Choose the Direction of Maximal Variance}

Suppose \(v\) is an eigenvector of \(S\), so
\[
Sv=\lambda v.
\]

Then
\[
v^TSv
=
v^T(\lambda v)
=
\lambda v^Tv.
\]

Because
\[
v^Tv=1,
\]
we get
\[
v^TSv=\lambda.
\]

Therefore, the variance along direction \(v\) is proportional to the eigenvalue \(\lambda\).

Thus, to maximize the projected variance, we choose the eigenvector with the largest eigenvalue.

This direction is called the first principal component.

\section*{Step 6: Extend to \(k\) Dimensions}

For a \(k\)-dimensional projection, we choose \(k\) orthonormal directions
\[
u_1,u_2,\dots,u_k.
\]

Collect them into the matrix
\[
U_k = [u_1,u_2,\dots,u_k]
\in \mathbb{R}^{p\times k}.
\]

The orthonormality condition is
\[
U_k^TU_k=I_k.
\]

The projected representation is
\[
y_i = U_k^Tz_i
=
U_k^T(x_i-m).
\]

Since each eigenvalue measures the variance preserved by its eigenvector, the best \(k\)-dimensional subspace is obtained by choosing the eigenvectors corresponding to the \(k\) largest eigenvalues:
\[
U_k = [u_1,u_2,\dots,u_k],
\]
where
\[
\lambda_1 \geq \lambda_2 \geq \cdots \geq \lambda_p.
\]

\section*{Step 7: Reconstruction and Approximation Error}

The projected coordinates are
\[
y_i = U_k^Tz_i.
\]

The reconstruction in the original space is
\[
\hat z_i = U_k y_i.
\]

Substituting \(y_i = U_k^Tz_i\), we get
\[
\hat z_i = U_kU_k^Tz_i.
\]

Thus, the reconstruction of the original point is
\[
\hat x_i = m + U_kU_k^T(x_i-m).
\]

The reconstruction error is
\[
z_i-\hat z_i
=
z_i-U_kU_k^Tz_i
=
(I-U_kU_k^T)z_i.
\]

Therefore, the total approximation error is
\[
\sum_{i=1}^n
\left\|
(I-U_kU_k^T)(x_i-m)
\right\|^2.
\]

Choosing the eigenvectors with the largest eigenvalues maximizes the retained variance.

Since the total variance is fixed, maximizing retained variance is equivalent to minimizing the remaining reconstruction error.

\section*{Conclusion}

PCA starts from the goal
\[
y_i = U_k^T(x_i-m),
\]
where \(y_i\) is a lower-dimensional representation of \(x_i\).

To preserve as much information as possible, PCA maximizes the variance of the projected data. This leads to the constrained optimization problem
\[
\max_v v^TSv
\quad
\text{subject to}
\quad
v^Tv=1.
\]

Using Lagrange multipliers, this becomes the eigenvalue problem
\[
Sv=\lambda v.
\]

Therefore, PCA chooses the eigenvectors of the scatter matrix with the largest eigenvalues.

These directions preserve the largest amount of variance and minimize the approximation error.

\[
\boxed{
\text{PCA reduces dimensionality by projecting data onto the directions of maximal variance.}
}
\]

\end{document}